\section{Experimentos}
\label{sec:experimentos}

Reserve de uma a duas páginas para a avaliação. Uma combinação eficaz costuma
ser uma tabela de métricas, um gráfico e um parágrafo de discussão que
interprete os números, e não apenas os repita.

\subsection{Metodologia}
\label{sec:experimentos:metodologia}

Descreva o conjunto de dados, os modelos ou algoritmos comparados, os
hiperparâmetros, as métricas, o hardware utilizado e tudo o que for necessário
para reproduzir os resultados.

\subsection{Resultados}
\label{sec:experimentos:resultados}

Apresente os resultados e analise-os de forma crítica. A Tabela~\ref{tab:resultados}
ilustra um formato possível para reportar métricas por configuração.

\begin{table}[H]
  \centering
  \caption{Resultados comparativos entre as configurações avaliadas.}
  \label{tab:resultados}
  \begin{tabular}{lcc}
    \toprule
    Configuração & Métrica A & Métrica B \\
    \midrule
    Linha de base   & 0{,}71 & 0{,}65 \\
    Proposta        & 0{,}88 & 0{,}82 \\
    \bottomrule
  \end{tabular}
\end{table}

Comente o que os números revelam, em que condições a proposta supera a linha de
base e quais limitações ainda persistem.
